\section{\perfrace{}: Benchmark Details}
\label{app:perfrace}

The \perfrace{} appendix is organized as follows: Section~\ref{app:perfrace-comparison} compares \perfrace{} with related benchmarks; Section~\ref{app:perfrace-characterization} characterizes its 100 tasks; Section~\ref{app:perfrace-curation-statistics} reports construction and curation statistics; Section~\ref{app:perfrace-pair-verification} describes the Pair Verification protocol and quality control; and Section~\ref{app:perfrace-evaluation} details the evaluation protocol and analyzes gains from longer optimization budgets.

\subsection{Benchmark Curation and Evaluation}
\label{app:perfrace-comparison}

Table~\ref{tab:perfrace-benchmark-comparison} compares \perfrace{} with
representative ultra-long-horizon agent benchmarks and code-optimization
benchmarks along task type, task construction, and evaluation procedure.
Existing ultra-long-horizon benchmarks largely assemble fixed suites of
hand-crafted or expert-curated tasks. In contrast, \perfrace{} combines live
GitHub discovery with agent--human Pair Verification to continually construct
validated optimization tasks from real-world projects and workloads.
Compared with existing code-optimization benchmarks, it does not derive tasks
from known performance patches and evaluates progress repeatedly rather than
only at task completion.

A workload specifies the inputs and operations used to exercise a project.
A microbenchmark measures a narrowly scoped part of that work, such as a
selected API call or operation. For example, SWE-fficiency includes a
performance workload that repeatedly times a NumPy \texttt{einsum}
call~\citep{ma2025swefficiency}.
An end-to-end workload instead measures a complete project-use workflow.
\perfrace{} adopts this broader measurement boundary, leaving agents to
profile the workflow and identify which components to optimize.
The distinction is therefore the scope of the measured computation:
targeted operations versus an end-to-end workflow, rather than whether
the benchmark supplies a workload.

\begin{table}[H]
  \centering
  \scriptsize
  \setlength{\tabcolsep}{3pt}
  \renewcommand{\arraystretch}{1.25}
  \caption{Comparison with related benchmarks. Microbenchmarks measure targeted
  operations (for example, selected API calls), whereas end-to-end workloads
  measure complete project-use workflows.}
  \label{tab:perfrace-benchmark-comparison}
  \begin{tabularx}{\textwidth}{@{}l
    >{\hsize=0.9\hsize\linewidth=\hsize\raggedright\arraybackslash\hspace{0pt}}X
    >{\hsize=1.3\hsize\linewidth=\hsize\raggedright\arraybackslash\hspace{0pt}}X
    >{\hsize=0.8\hsize\linewidth=\hsize\raggedright\arraybackslash\hspace{0pt}}X@{}}
    \toprule
    Benchmark & Task type & Task construction & Evaluation \\
    \midrule
    \perfrace{} &
    Repository optimization on end-to-end workloads &
    Live GitHub discovery followed by agent--human Pair Verification &
    Iterative correctness and performance feedback from an external Judge \\
    AutoLab~\mbox{\citep{xu2026autolab}} &
    Diverse metric-driven optimization tasks &
    Expert-curated optimization tasks &
    Local evaluation during the task; held-out final evaluation \\
    EdgeBench~\mbox{\citep{zhu2026edgebench}} &
    Cross-domain iterative problem-solving tasks &
    Expert-curated day-scale tasks &
    Local inner loop; repeated submission-gated Judge feedback \\
    FrontierSWE~\mbox{\citep{chu2026frontierswe}} &
    Ultra-long-horizon software engineering and research tasks &
    Real-world problems curated with academic and industry experts &
    Final graded submission after up to 20 hours \\
    Terminal-Bench~\mbox{\citep{merrill2026terminalbench}} &
    General terminal tasks &
    Crowdsourced tasks with multi-stage human review &
    Final-only test-based evaluation \\
    GSO~\mbox{\citep{shetty2025gso}} &
    Microbenchmark-based repository optimization &
    Repository commit histories and generated performance tests &
    Final-only microbenchmark-based evaluation \\
    SWE-Perf~\mbox{\citep{he2025sweperf}} &
    Microbenchmark-based repository optimization &
    Historical performance PRs and selected repository tests &
    Final-only microbenchmark-based evaluation \\
    SWE-fficiency~\mbox{\citep{ma2025swefficiency}} &
    Microbenchmark-based repository optimization &
    Performance PRs and manually reconstructed workloads &
    Final-only microbenchmark-based evaluation \\
    FormulaCode~\mbox{\citep{sehgal2026formulacode}} &
    Microbenchmark-based repository optimization &
    Continuously ingested performance PRs and ASV workloads &
    Final-only microbenchmark-based evaluation \\
    \bottomrule
  \end{tabularx}
\end{table}


\subsection{Detailed Benchmark Composition}
\label{app:perfrace-characterization}

We characterize the 100 \perfrace{} tasks by project purpose, primary programming language, license, packaged source size, file count, and fixed workload.

\textbf{Projects.}
The projects represented in \perfrace{} span 16 primary programming languages:
Rust (25 tasks), Python (18), TypeScript (15), C++ (11), C (6), C\# (6), Julia
(5), Go (4), Java (3), Scala (1), Lua (1), OpenSCAD (1), Solidity (1), MATLAB
(1), JavaScript (1), and Zig (1).
The packaged project sources range from 42 to 15,869 files and from 471.49 KiB
to 689.08 MiB, with medians of 886 files and 14.56 MiB.

\begin{figure}[H]
  \centering
  \includegraphics[width=\textwidth]{figures/perfrace-projects.pdf}
  \caption{Distribution of primary programming languages, packaged source sizes, and packaged source file counts across \perfrace{} projects.}
  \label{fig:perfrace-projects}
\end{figure}

\textbf{Task catalog.}
Table~\ref{tab:perfrace-task-catalog} lists each task's project, purpose,
license, and fixed workload.

{\scriptsize
\setlength{\tabcolsep}{3pt}
\renewcommand{\arraystretch}{1.12}
\begin{longtable}{@{}>{\raggedright\arraybackslash}p{0.22\textwidth}
  >{\raggedright\arraybackslash}p{0.23\textwidth}
  >{\raggedright\arraybackslash}p{0.13\textwidth}
  >{\raggedright\arraybackslash}p{0.36\textwidth}@{}}
  \caption{Complete catalog of the 100 \perfrace{} tasks.}
  \label{tab:perfrace-task-catalog}\\
  \toprule
  Task & Project and purpose & License & Fixed workload \\
  \midrule
  \endfirsthead
  \multicolumn{4}{c}{\tablename~\thetable{} (continued)} \\
  \toprule
  Task & Project and purpose & License & Fixed workload \\
  \midrule
  \endhead
  \midrule
  \multicolumn{4}{r}{Continued on next page} \\
  \endfoot
  \bottomrule
  \endlastfoot
  \path{aiidateam-aiida-core--battery-archive-import} & \path{aiidateam/aiida-core}\newline Scientific workflow and provenance platform & MIT & AiiDA import of a public battery-cycling provenance archive \\
  \path{anyrender--openmoji} & \path{DioxusLabs/anyrender}\newline Rust 2D drawing abstraction & Apache-2.0 & AnyRender batch rasterization of the OpenMoji SVG corpus \\
  \path{apalache--zcash-sync} & \path{apalache-mc/apalache}\newline Symbolic model checker for TLA+ and Quint & Apache-2.0 & Apalache's bounded safety check for the Zcash v2 block-download scheduler \\
  \path{arrow-rs--variant-wikidata} & \path{apache/arrow-rs}\newline Rust implementation of Apache Arrow & Apache-2.0 & apache/arrow-rs while encoding Wikidata entities as a Variant column \\
  \path{bevy-voxel-world--terrain-streaming} & \path{splashdust/bevy_voxel_world}\newline Voxel-world plugin for Bevy & Apache-2.0 & splashdust/bevy\_voxel\_world terrain generation and chunk meshing \\
  \path{bm25s--trec-covid-index} & \path{xhluca/bm25s}\newline BM25 search library & MIT & bm25s indexing and persistence for the TREC-COVID research collection \\
  \path{boyter-scc--linux-v6-19-json-scan} & \path{boyter/scc}\newline Source-code counting and complexity tool & MIT & SCC whole-repository JSON scan time on Linux v6.19 \\
  \path{calldiff--pi-review} & \path{tanishqkancharla/calldiff}\newline Function call-stack diff tool & MIT & Calldiff JSON call-stack diffs for a real Pi coding-agent review \\
  \path{cbcoutinho-nextcloud-mcp-server--ohr-pdf-ingest} & \path{cbcoutinho/nextcloud-mcp-server}\newline MCP server for Nextcloud & AGPL-3.0 & Nextcloud MCP Server's born-digital PDF ingest and page-aware chunking throughput \\
  \path{ccc--compile-speed} & \path{anthropics/claudes-c-compiler}\newline Dependency-free C compiler & CC0-1.0 & CCC compile-to-object latency on the SQLite 3.46.0 amalgamation \\
  \path{chaste--ttp06-sheet} & \path{Chaste/Chaste}\newline Biomedical simulation framework & BSD-3-Clause & Chaste cardiac simulation for a TTP06 ventricular sheet propagation \\
  \path{cmtraceopen--hadoop-logs} & \path{adamgell/cmtraceopen}\newline Cross-platform Windows log viewer & MIT & CMTrace Open while loading a real Hadoop/YARN log session \\
  \path{cocoon--oci-import} & \path{cocoonstack/cocoon}\newline MicroVM engine & MIT & local import of a supported Cocoon Ubuntu VM OCI image \\
  \path{codediff--neovim-release} & \path{esmuellert/codediff.nvim}\newline Neovim code-review and diff plugin & MIT & CodeDiff C diff latency for a real Neovim release review \\
  \path{commitizen--check} & \path{commitizen-tools/commitizen}\newline Commit and changelog automation tool & MIT & Commitizen changelog generation for Angular's complete Git history \\
  \path{core--blob-upload} & \path{calimero-network/core}\newline Peer-to-peer application platform & Apache-2.0 & Calimero content-addressed blob ingestion for a real OpenStreetMap file upload \\
  \path{cpp2rust--woff2-compress} & \path{Cpp2Rust/cpp2rust}\newline C++-to-safe-Rust translator & MIT & Cpp2Rust's safe Rust output for public-font WOFF2 compression \\
  \path{deepinv--kohler-kernels} & \path{deepinv/deepinv}\newline Deep-learning inverse-problems library & BSD-3-Clause & DeepInverse spatially varying blur-kernel estimation on a real Kohler camera-motion-blurred image \\
  \path{dftd4--lysozyme-gradient} & \path{dftd4/dftd4}\newline Atomic-charge-dependent dispersion correction & LGPL-3.0 & DFT-D4 energy, analytic gradient, and virial evaluation for a real lysozyme structure \\
  \path{dither-guy--photo-batch} & \path{manoelpiovesan/dither-guy}\newline Image dithering application & GPL-3.0 & Dither Guy's full-resolution photo batch export \\
  \path{doc-detective--redpanda-transforms} & \path{doc-detective/doc-detective}\newline Documentation testing framework & AGPL-3.0 & Doc Detective dry-run planning for Redpanda transform documentation \\
  \path{drasi-core--gtfs-delay-replay} & \path{drasi-project/drasi-core}\newline Continuous-query engine & Apache-2.0 & Drasi Core's continuous Cypher evaluation for replayed GTFS-Realtime vehicle geofence updates \\
  \path{duckdb--taxi-ingest} & \path{Giorgi/DuckDB.NET}\newline .NET bindings for DuckDB & MIT & DuckDB.NET ingestion of NYC Yellow Taxi trip records \\
  \path{dusty--fpindex-ingest} & \path{lalinsky/dusty}\newline Zig HTTP client/server library & MIT & Dusty's HTTP batch-ingestion path for real AcoustID fpindex updates \\
  \path{embree--chess-render} & \path{RenderKit/embree}\newline Ray-tracing kernels & Apache-2.0 & Embree triangle traversal for a public chess-scene visibility render \\
  \path{endee--esci} & \path{endee-io/endee}\newline Vector database & AGPL-3.0 & Endee HTTP dense search for a replay of Amazon Shopping Queries product searches \\
  \path{enzyme--mnist} & \path{EnzymeAD/Enzyme.jl}\newline Julia automatic-differentiation library & MIT & Enzyme.jl for Flux's CPU LeNet MNIST training epoch \\
  \path{eole--gpt2-wikitext-generate} & \path{eole-nlp/eole}\newline Language-modeling toolkit & MIT & EOLE CPU GPT-2 batch text generation on WikiText-103 prompts \\
  \path{eucalypt--config-render-speed} & \path{curvelogic/eucalypt}\newline Configuration language and interpreter & MIT OR Apache-2.0 & the eucalypt interpreter on a config-generation workload \\
  \path{eucalypt--eval-speed} & \path{curvelogic/eucalypt}\newline Configuration language and interpreter & MIT OR Apache-2.0 & the eucalypt interpreter on an Advent of Code workload \\
  \path{facet--venta-guild} & \path{Tim-Maes/Facet}\newline .NET DTO and mapping source generator & MIT & Facet package and source-generator compilation time for venta's Guild service \\
  \path{flexmeasures--smard-price-import} & \path{FlexMeasures/flexmeasures}\newline Energy-management platform & Apache-2.0 & FlexMeasures import of a full year of public SMARD day-ahead prices \\
  \path{fluidsynth-fluidsynth--bagatelles-midi-render} & \path{FluidSynth/fluidsynth}\newline SoundFont software synthesizer & LGPL-2.1 & FluidSynth offline rendering of a complete CC0 MIDI performance to WAV \\
  \path{geoparquet-io--hilbert-sort} & \path{geoparquet/geoparquet-io}\newline GeoParquet tooling & Apache-2.0 & GeoParquet Hilbert ordering of an Overture Buildings shard \\
  \path{gridflock--drawer-stl} & \path{yawkat/GridFlock}\newline Gridfinity baseplate generator & MIT / CC-BY-4.0 & GridFlock's segmented Gridfinity drawer-baseplate STL export \\
  \path{ifcopenshell--gni-bimproject-obj-export} & \path{IfcOpenShell/IfcOpenShell}\newline IFC library and geometry engine & LGPL-3.0 & IfcOpenShell IfcConvert OBJ export for a lighter GNI BIM Project architectural IFC delivery \\
  \path{intro-skipper--credits} & \path{intro-skipper/intro-skipper}\newline Jellyfin intro and credits detector & GPL-3.0 & Intro Skipper's movie credits black-frame evidence extraction \\
  \path{jagua--nesting} & \path{JeroenGar/jagua-rs}\newline 2D collision engine for nesting & MPL-2.0 & jagua-rs collision collection replayed from a downstream nesting solver \\
  \path{javaevolved--site-build} & \path{javaevolved/javaevolved.github.io}\newline Java developer documentation site & MIT & the full localized GitHub Pages build for java.evolved \\
  \path{jsonnet--kubernetes-mixin-speed} & \path{google/jsonnet}\newline Data templating language & Apache-2.0 & the Jsonnet C++ interpreter on a Kubernetes monitoring configuration workload \\
  \path{juliaio-zarr-jl--oceananigans-timeseries-output} & \path{JuliaIO/Zarr.jl}\newline Julia implementation of Zarr array storage & MIT & Zarr.jl directory-store output for a real Oceananigans tracer time series \\
  \path{julialang-julia--trixi-jeans-instability} & \path{JuliaLang/julia}\newline Julia programming language & MIT & Julia runtime execution for a published Trixi Jeans gravitational-instability simulation \\
  \path{kitops--minilm-pack} & \path{kitops-ml/kitops}\newline OCI packaging tool for ML artifacts & Apache-2.0 & KitOps local ModelKit packaging for a complete all-MiniLM-L6-v2 release snapshot \\
  \path{kiwifs--kubernetes-reindex} & \path{kiwifs/kiwifs}\newline Markdown filesystem for agents and teams & BUSL-1.1 & KiwiFS SQLite FTS5 reindex for Kubernetes documentation \\
  \path{kms--batch-encrypt} & \path{Cosmian/kms}\newline Key-management system & BUSL-1.1 & Eviden KMS KMIP BulkData AES-GCM encryption for concurrent application-record batches \\
  \path{kubb--esi} & \path{kubb-labs/kubb}\newline OpenAPI code-generation framework & MIT & Kubb regeneration of JitaSpace's EVE Online ESI client \\
  \path{latexsnipper--formula-batch} & \path{SakuraMathcraft/LaTeXSnipper}\newline Formula OCR and editing tool & GPL-3.0 & resident CPU formula OCR for a public OpenStax textbook screenshot batch \\
  \path{libcachesim--twitter-lru} & \path{1a1a11a/libCacheSim}\newline Cache simulation library & Apache-2.0 & Four-capacity LRU cache-size sweep over a production Twitter request trace \\
  \path{lsp-tree-sitter--termux-check} & \path{neomutt/lsp-tree-sitter}\newline Tree-sitter language-server library & GPL-3.0 & Tree-sitter validation of a real Termux package-recipe corpus \\
  \path{mai-gen-videob50--score-images} & \path{Nick-bit233/mai-gen-videob50}\newline Maimai score-video generation tool & MIT & batch maimai score-image rendering for a real 50-record archive \\
  \path{meshmonitor--coverage-grid} & \path{Yeraze/meshmonitor}\newline Mesh-node monitoring web tool & BSD-3-Clause & MeshMonitor's SQLite coverage-grid batch \\
  \path{metamorpho--reallocate} & \path{morpho-org/metamorpho}\newline Tokenized lending-vault protocol & GPL-2.0 & Gas use for a historical MetaMorpho market reallocation \\
  \path{mini-qr--library-batch} & \path{lyqht/mini-qr}\newline QR-code generator and scanner & GPL-3.0 & Mini QR's batch SVG export for a public library directory \\
  \path{mirage--kubernetes-triage} & \path{strukto-ai/mirage}\newline Virtual terminal for coding agents & Apache-2.0 & Mirage recursive source search for Kubernetes dual-stack readiness-probe triage \\
  \path{modflow6--neversink} & \path{MODFLOW-ORG/modflow6}\newline Groundwater-flow simulator & CC0-1.0 & Complete steady-state simulation of the USGS Neversink--Rondout groundwater model \\
  \path{monorepo--crypto-build} & \path{commonwarexyz/monorepo}\newline Distributed-systems primitives & Apache-2.0 & the release build latency of commonware-cryptography \\
  \path{mrst--norne-import} & \path{SINTEF-AppliedCompSci/MRST}\newline Reservoir-simulation toolbox & GPL-3.0 & MRST Norne field-model import and geometry construction \\
  \path{msdialworkbench--lcms-alignment} & \path{systemsomicslab/MsdialWorkbench}\newline Mass-spectrometry analysis workbench & LGPL-3.0 & MS-DIAL Console LC-MS peak processing and alignment for a real negative-ion DDA batch \\
  \path{mssql-extension--taxi-copy} & \path{hugr-lab/mssql-extension}\newline DuckDB extension for Microsoft SQL Server & MIT & BCP bulk import of NYC taxi records through the DuckDB MSSQL extension \\
  \path{n8n-mcp--tool-batch-speed} & \path{czlonkowski/n8n-mcp}\newline MCP server for n8n workflow construction & MIT & n8n-mcp on a fixed MCP tool-call batch \\
  \path{nerdbank-messagepack--earthquake-fanout} & \path{AArnott/Nerdbank.MessagePack}\newline .NET MessagePack library & MIT & SignalR MessagePack earthquake snapshot fan-out \\
  \path{ojp--taxi-export} & \path{Open-J-Proxy/ojp}\newline Database proxy and JDBC driver & Apache-2.0 & OJP JDBC export of public NYC TLC taxi records \\
  \path{omnigraph--metabolic-trials-refresh} & \path{ModernRelay/omnigraph}\newline Lakehouse-native graph engine & MIT & OmniGraph's paginated metabolic clinical-trial refresh and lifecycle dashboard \\
  \path{onnxscript--bert-optimize} & \path{microsoft/onnxscript}\newline ONNX model authoring and optimization library & MIT & Graph optimization of a BERT-SQuAD FP32 and INT8 model bundle \\
  \path{oxidize-pdf--rag-chunking} & \path{bzsanti/oxidizePdf}\newline Rust PDF library for RAG & MIT & oxidize-pdf RAG chunking over a real CC-BY PDF corpus \\
  \path{pdal--dtm} & \path{PDAL/PDAL}\newline Point-cloud processing library & BSD-3-Clause & Ground classification and 1 m terrain-model generation from USGS 3DEP LiDAR \\
  \path{petsc--ocean-age-solve} & \path{JuliaParallel/PETSc.jl}\newline Julia wrappers for PETSc & MIT & PETSc.jl GMRES+ILU for a public OCCA ocean-age transport solve \\
  \path{physical-ai-toolchain--action-smoothness} & \path{microsoft/physical-ai-toolchain}\newline Physical-AI dataset analysis toolkit & MIT & multi-dataset LeRobot action-smoothness analysis for real ALOHA manipulation recordings \\
  \path{plonky3--ethereum-calldata-keccak-proof} & \path{Plonky3/Plonky3}\newline Polynomial IOP toolkit & Apache-2.0 & Plonky3 Ethereum calldata Keccak trace generation and STARK proving \\
  \path{polkadot-api--descriptor-generation} & \path{polkadot-api/polkadot-api}\newline Modular Polkadot API libraries & MIT & PAPI's multi-chain TypeScript descriptor generation \\
  \path{pytrendy--dow-index} & \path{RussellSB/pytrendy}\newline Trend-detection library & MIT & PyTrendy trend detection for a real daily Dow Jones Industrial Average price series \\
  \path{qrisp--modmul-compile} & \path{eclipse-qrisp/Qrisp}\newline Quantum-algorithm development framework & EPL-2.0 & Qrisp compilation of controlled modular multiplication \\
  \path{react-native-code-push--mattermost-release} & \path{Soomgo-Mobile/react-native-code-push}\newline React Native OTA update client & MIT & CodePush Android OTA releases for Mattermost Mobile \\
  \path{reactant--fashion-mnist-train} & \path{EnzymeAD/Reactant.jl}\newline MLIR/XLA compiler for Julia & MIT & Reactant CPU execution of a real Lux Fashion-MNIST training epoch \\
  \path{rspress--react-native-paper} & \path{web-infra-dev/rspress}\newline Static-site generator & MIT & Rspress production builds of the React Native Paper documentation site \\
  \path{rust-clippy--cargo-package-lint} & \path{rust-lang/rust-clippy}\newline Rust linter & Apache-2.0 & Clippy lint time for Cargo's main production package \\
  \path{s3s--kubernetes-sync} & \path{s3s-project/s3s}\newline S3 service adapter & Apache-2.0 & s3s-fs recursive S3 synchronization of Kubernetes pkg source \\
  \path{senpi--image-batch} & \path{code-yeongyu/senpi}\newline Coding-agent runtime & MIT & Senpi image attachment preprocessing for a real screenshot batch \\
  \path{snarkos--cdn-replay} & \path{ProvableHQ/snarkOS}\newline Decentralized operating system for ZK applications & Apache-2.0 & snarkOS CDN block bundle replay over 150 fixed Aleo Mainnet blocks \\
  \path{snarkvm--mainnet-execution-verify} & \path{ProvableHQ/snarkVM}\newline Zero-knowledge virtual machine & Apache-2.0 & snarkVM verification of a batch of accepted Aleo mainnet executions \\
  \path{solvespace-solvespace--printer-assembly-mesh-export} & \path{solvespace/solvespace}\newline Parametric 2D/3D CAD & GPL-3.0 & SolveSpace CLI mesh export for a complete Cartesian 3D-printer assembly \\
  \path{spectral--api-lint} & \path{stoplightio/spectral}\newline OpenAPI, JSON, and YAML linter & Apache-2.0 & Spectral linting of a production OpenAPI batch \\
  \path{sqruff--glam-etl} & \path{quarylabs/sqruff}\newline SQL formatter and linter & Apache-2.0 & Sqruff directory linting for Mozilla GLAM ETL SQL \\
  \path{strawberry-graphql-strawberry--graphique-tlc-weekly-columns} & \path{strawberry-graphql/strawberry}\newline Python GraphQL library & MIT & Strawberry GraphQL execution for Graphique's real TLC weekly Parquet column query \\
  \path{suitesparse--livejournal-pagerank} & \path{DrTimothyAldenDavis/SuiteSparse}\newline Sparse-matrix and graph computation suite & Apache-2.0 / BSD-2-Clause (used components) & Complete PageRank computation on the LiveJournal social graph using GraphBLAS and LAGraph \\
  \path{sum--text-attest} & \path{OtotaO/SUM}\newline Cross-runtime text-attestation library & Apache-2.0 & SUM while batch-attesting a corpus of real Wikipedia prose \\
  \path{sundials--cantera-ignition} & \path{llnl/sundials}\newline Nonlinear and differential-equation solver suite & BSD-3-Clause & SUNDIALS CVODES for a downstream Cantera methane ignition parameter sweep \\
  \path{svf-tools-svf--nginx-andersen-wpa} & \path{SVF-tools/SVF}\newline Static value-flow analysis framework & AGPL-3.0-or-later & SVF Andersen whole-program pointer analysis on NGINX reverse-proxy bitcode \\
  \path{synalinks--kubernetes-tasks} & \path{SynaLinks/synalinks}\newline Neuro-symbolic language-model framework & Apache-2.0 & Synalinks Markdown knowledge-base ingestion for Kubernetes task documentation \\
  \path{tfhe-rs--erc7984-transfer} & \path{zama-ai/tfhe-rs}\newline Fully homomorphic encryption library & BSD-3-Clause-Clear & TFHE-rs CPU execution for confidential ERC-7984 token transfers \\
  \path{tileserver-gl--country-atlas} & \path{maptiler/tileserver-gl}\newline Vector and raster map-tile server & BSD-2-Clause & TileServer GL static PNG rendering for a full Natural Earth country atlas \\
  \path{trixi-framework-trixi-jl--self-gravitating-sedov-amr} & \path{trixi-framework/Trixi.jl}\newline Adaptive numerical-simulation framework & MIT & a published self-gravitating Sedov blast-wave simulation with adaptive mesh refinement \\
  \path{typst--napkin} & \path{Myriad-Dreamin/typst.ts}\newline Typst runtime for JavaScript & Apache-2.0 & Typst.ts rendering in the Napkin Community production static-site build \\
  \path{vmecpp--precise-qh} & \path{proximafusion/vmecpp}\newline Magnetohydrodynamic equilibrium solver & MIT & VMEC++ fixed-boundary precise-QH equilibrium solves \\
  \path{wiredtiger--pageviews} & \path{wiredtiger/wiredtiger}\newline Embedded storage engine & GPL-2.0 OR GPL-3.0 & WiredTiger's durable import of a real Wikimedia pageview hour \\
  \path{xet-core--bert-download} & \path{huggingface/xet-core}\newline Content-addressed transfer and reconstruction engine & Apache-2.0 & local Xet reconstruction of a production BERT model artifact \\
  \path{yamato-security-hayabusa--external-attack-evtx-timeline} & \path{Yamato-Security/hayabusa}\newline Windows event-log forensics tool & AGPL-3.0 & Hayabusa bulk external attack-EVTX JSONL timeline generation \\
  \path{youtrackdb--osm-roads} & \path{JetBrains/youtrackdb}\newline Object-oriented graph database & Apache-2.0 & embedded YouTrackDB import of a real OpenStreetMap road network \\
  \path{zenbu--release-source} & \path{zenbu-labs/zenbu.js}\newline JavaScript application framework & MIT & Zenbu.js staging of the Zenbu desktop application's release source \\
  \path{zio--zrk-typecheck} & \path{lalinsky/zio}\newline Zig asynchronous-I/O framework & MIT & ZIO's downstream zrk ReleaseFast type-check latency \\
\end{longtable}
}


\subsection{Construction and Curation Statistics}
\label{app:perfrace-curation-statistics}

This section describes the implementation of Discovery and Filter and reports the candidate counts and curation effort at each construction stage.

\textbf{Discovery implementation.}
We query the ClickHouse \texttt{github\_events} table for opened pull requests and branch pushes by a configured list of coding-agent accounts, then group matching events by repository.
For each repository, we retain the latest matching pull-request or branch URL and use the GitHub GraphQL API to record its canonical identity, default-branch revision, stars, primary language, license, disk usage, and repository status.

The resulting records are deduplicated by GitHub repository, and projects are excluded when their repository, default branch, or recorded source revision cannot be retrieved.
Each step writes its results incrementally, allowing interrupted collection and metadata retrieval to resume without repeating completed work.

From GitHub activity recorded between May and August 2026, we identified 118,517 repositories, retrieved metadata for 104,428, and retained 103,475 project candidates after these checks.

\begin{table}[H]
  \centering
  \small
  \setlength{\tabcolsep}{10pt}
  \renewcommand{\arraystretch}{1.15}
  \caption{GitHub activity signals used for Discovery.}
  \label{tab:perfrace-discovery-statistics}
  \begin{tabular}{@{}lr@{}}
    \toprule
    Signal & Repositories \\
    \midrule
    Branch push by a coding-agent account & 112,127 \\
    Pull request opened by a coding-agent account & 6,390 \\
    \bottomrule
  \end{tabular}
\end{table}


\textbf{Filter implementation.}
We first apply deterministic checks to repository metadata and the recorded Git tree.
These checks require an accessible source revision, a license permitting modification and redistribution, 10--4,999 GitHub stars, and a tracked Git tree of at most 500 MiB and 25,000 entries.

Projects that pass these checks are collected at the recorded revision with their README, setup documentation, package manifests, license files, and representative source paths.
We then use GPT-5.6-Luna to assess the project's setup path, evaluable process, runtime fit, and optimization potential from this source snapshot using a structured response schema.
The prompt is provided in Appendix~\ref{app:perfrace-filter-prompt}.

We require each structured response to address the following criteria.
\begin{list}{$\square$}{
  \setlength{\leftmargin}{1.6em}
  \setlength{\labelwidth}{1em}
  \setlength{\labelsep}{0.4em}
  \setlength{\topsep}{2pt}
  \setlength{\itemsep}{1pt}
  \setlength{\parsep}{0pt}
  \setlength{\partopsep}{0pt}
}
  \item \textbf{Setup path.} Repository files support concrete commands for installing dependencies and building or running the project.
  \item \textbf{Evaluable process.} The project supports a finite process with a fixed input, observable output, performance metric, and correctness check.
  \item \textbf{Runtime fit.} The process can run in a Linux x86-64 container on commodity CPUs without mandatory special hardware.
  \item \textbf{Optimization potential.} The project supports a substantial local computation whose cost is expected to increase with input size or repeated execution.
\end{list}

Using these assessments, we selected 239 high-potential project candidates for Pair Verification.
Table~\ref{tab:perfrace-filter-statistics} reports the first check that excludes each project from the final Pair Verification selection.

\begin{table}[H]
  \centering
  \small
  \setlength{\tabcolsep}{10pt}
  \renewcommand{\arraystretch}{1.15}
  \caption{First failed check during Filter and the final selection result.}
  \label{tab:perfrace-filter-statistics}
  \begin{tabular}{@{}lr@{}}
    \toprule
    Result & Projects \\
    \midrule
    Fewer than 10 GitHub stars & 93,992 \\
    Accessibility & 1,134 \\
    Setup path & 212 \\
    Runtime fit & 359 \\
    Optimization potential & 4,462 \\
    Evaluable process & 3,066 \\
    Assessment error & 11 \\
    \midrule
    Selected for Pair Verification & 239 \\
    \bottomrule
  \end{tabular}
\end{table}


\textbf{Curation effort.}
To quantify the effort required to turn selected projects into accepted tasks, Figure~\ref{fig:perfrace-curation-effort} reports the agent cost, curation time, and number of agent--human interactions for the 100 accepted tasks.
The median task required 4.5 hours, \$8.83 in agent cost, and eight agent--human interactions, compared with averages of 9.1 hours, \$11.35, and 12 interactions.

\begin{figure}[H]
  \centering
  \includegraphics[width=\textwidth]{figures/perfrace-curation-effort.pdf}
  \caption{Distribution of curation time, agent cost, and agent--human interactions across the 100 accepted tasks.}
  \label{fig:perfrace-curation-effort}
\end{figure}

\subsection{Pair Verification}
\label{app:perfrace-pair-verification}

\perfrace{} uses Pair Verification to ensure that each task represents meaningful project use, produces reproducible measurements, and rejects shortcut solutions; this section describes the agent--human protocol and analyzes the resulting curation trajectories.

\subsubsection{Protocol}

Pair Verification combines agent-led task construction with human review to turn each project candidate into an executable task and assess its realism, reliability, and integrity.
The construction skill is provided in Appendix~\ref{app:perfrace-pair-verification-skill}.
The agent develops the task in stages and pauses for a reviewer decision after completing the evidence required for each quality criterion.
\begin{figure}[H]
  \centering
  \includegraphics[width=\textwidth]{figures/perfrace-pair-verification.pdf}
  \caption{Pair Verification alternates between agent-led task construction and human review at the Realism, Reliability, and Integrity gates.}
  \label{fig:perfrace-pair-verification}
\end{figure}

\textbf{Explore and Design.}
The agent first studies the project's documentation, examples, code paths, and public usage evidence to understand its purpose, users, and performance-relevant operations.
The agent searches public release pages for source code, models, and assets, scientific or government data portals, downstream repositories, production archives, and real user exports to identify candidate workload inputs.
For each input, the agent records its download location, version, and applicable terms, and accepts it only when an explicit license, public-domain designation, or written authorization permits redistribution with the benchmark.

It uses these materials to identify candidate operations, but grounds the workload and scale in a real user or downstream application rather than repository-native tests, benchmarks, examples, or fixtures.
The agent compares plausible workload candidates, then traces the selected operation from its user-facing entry point through the relevant project code and demonstrates a small input-to-output execution on the unmodified revision.
During this run, it records dependency installation, project build, input preparation, and operation time to check that setup does not dominate the measurement and to estimate whether complete verification can finish within the 1,200-second budget.
The agent repeats the operation on identical inputs to detect run-to-run flakiness, including sporadic failures or inconsistent correctness-relevant outputs, and to identify any dependence of runtime on random state, network access, or external services.
A candidate is revised or rejected if its stochastic execution remains too variable for reproducible performance measurement after fixing available random seeds and aggregating repeated runs.

It then specifies the exact operation, workload and scale available to the agent, held-out correctness cases, observable correctness utility, performance metric, and measurement boundary.
Before requesting Realism review, the agent summarizes the project's purpose and users, the selected execution path, the basis for the input and scale, and the alternatives it considered and rejected.
The reviewer approves Realism only when the operation, input, and scale represent meaningful project use, the metric measures the stated utility, and the workload inputs can be redistributed as part of the benchmark.

\textbf{Build and Calibrate.}
After Realism approval, the agent derives the installation, build, and run commands from upstream documentation, CI configuration, and package manifests, and verifies them on the pinned source revision.
The agent implements the verifier, the task's test program that measures the performance metric on a separate copy of the workload and checks additional held-out cases for correctness.
It then constructs the Harbor task files with \texttt{instruction.md}, \texttt{task.toml}, an \texttt{environment/} directory containing the editable project and workload, and a separate \texttt{tests/} directory containing the verifier code, workload copy, held-out cases, and expected outputs.
The agent pins the project revision, base image, language toolchain, dependencies, workload assets, and expected outputs so that the image can be rebuilt and the task rerun without downloading project source or workload inputs.
The task instruction tells the agent how to build and run the project, what it may edit, which behavior it must preserve, and which performance metric to improve.

The agent then builds a fresh image and runs the verifier on the unmodified project before calibrating performance measurements.
The agent selects the smallest resource profile that completes the workload without CPU, memory, or storage pressure and sets the verifier timeout within 1,200 seconds.
The agent uses pilot and repeated runs to choose the workload scale, warmups, repetitions, and aggregation, requiring the final measurements to have a coefficient of variation below 10\% while keeping build and setup outside the timed operation.
The agent checks that the timing boundary covers the complete operation requested in the task while excluding one-time environment setup and verifier preparation that are unrelated to the optimization objective.

For the \(n\) timed measurements retained after warmup, \(x_1,\ldots,x_n\), we compute
\[
\mathrm{CV}=\frac{s}{\bar{x}},
\qquad
\bar{x}=\frac{1}{n}\sum_{i=1}^{n}x_i,
\qquad
s=\sqrt{\frac{1}{n-1}\sum_{i=1}^{n}(x_i-\bar{x})^2}.
\]

After fixing the workload, correctness checks, resources, and measurement protocol, the agent builds the final image and evaluates the unmodified baseline and one optimization candidate under the same conditions.
The reviewer examines the fresh-build correctness results, repeated measurements and CVs, verification time, resource allocation, and candidate trajectory.
The reviewer approves Reliability only when fresh builds and repeated runs return the same correctness verdict and stable performance measurements with the chosen warmups, repetitions, aggregation, and measurement boundary.

\textbf{Validate Integrity.}
To identify reward-hacking paths that could receive a valid score without completing the task, the agent tests deliberately incorrect, partial, hard-coded, and work-skipping solutions and checks that workload inputs, expected outputs, and verifier-owned data cannot be replaced by agent-controlled files.
The agent inspects the evaluation image for leaked expected outputs, verifier code, or construction notes that could reveal the hidden checks.
For example, in \texttt{rust-clippy--cargo-package-lint}, removing Cargo dependency tracking made later runs reuse a cached lint result, which the verifier detects by touching \texttt{clippy-driver}, running the workload twice against the same warm target, and requiring each Cargo JSON trace to report a non-fresh compiler artifact for \texttt{cargo}.

The reviewer approves Integrity only when every requirement in the task instruction is checked, every verifier check follows from the declared task, and none of the tested shortcut solutions can receive a valid score.

\textbf{Finalize.}
After all three reviews, the reviewer chooses \textsc{accept}, \textsc{revise}, or \textsc{reject}; accepted tasks pass final compatibility checks before publication as a task package comprising an immutable Docker image and task artifact.

\subsubsection{Quality Control}

We analyze the recorded agent--human Pair Verification trajectories for evidence of the Realism, Reliability, and Integrity of the resulting tasks.
Of the 239 candidates, 132 were rejected during Explore and Design, six during Build and Calibrate, none during Validate Integrity, and one during Finalize, leaving 100 accepted tasks.
\begin{figure}[H]
  \centering
  \includegraphics[width=0.78\textwidth]{figures/perfrace-stage-rejections.pdf}
  \caption{Candidates rejected at each Pair Verification stage, with candidates lacking a recorded transition into Build \& Calibrate assigned to Explore \& Design.}
  \label{fig:perfrace-stage-rejections}
\end{figure}
\textbf{Realism.}
Realism checks cover the workload operation and scale, the source and redistribution terms of its inputs, and the alignment between the declared utility and performance metric.
Table~\ref{tab:perfrace-workload-input-sources} groups the workload inputs into released source code, models, or assets (26 tasks), public scientific or government datasets (21), real downstream applications (21), production traces or archives (16), real user exports (6), and other independently grounded inputs (10).
\begin{table}[H]
  \centering
  \small
  \setlength{\tabcolsep}{6pt}
  \renewcommand{\arraystretch}{1.15}
  \caption{Sources of workload inputs across the 100 accepted tasks.}
  \label{tab:perfrace-workload-input-sources}
  \begin{tabularx}{\textwidth}{@{}>{\raggedright\arraybackslash}p{0.32\textwidth}
                                      >{\raggedright\arraybackslash}X
                                      r@{}}
    \toprule
    Input source & Examples & Tasks \\
    \midrule
    Released source code, model, or asset
      & Linux source tree, OpenMoji assets, released ML models & 26 \\
    Public scientific or government dataset
      & TREC-COVID, NYC taxi records, LC-MS data & 21 \\
    Real downstream application
      & React Native Paper, Headlamp, Mattermost Mobile & 21 \\
    Production trace or archive
      & Hadoop logs, GTFS delay feeds, GitHub events & 16 \\
    Real user export
      & MIDI collection, BIM project, printer assembly CAD & 6 \\
    Other independently grounded input
      & Wikidata variants, Dow index, OSM roads & 10 \\
    \bottomrule
  \end{tabularx}
\end{table}


\textbf{Reliability.}
All 99 wall-clock tasks have a final unmodified-project CV below 10\%; the remaining task, \texttt{metamorpho--reallocate}, measures transaction gas, the deterministic execution cost assigned by the Ethereum Virtual Machine to a fixed smart-contract transaction, and returns exactly 874,508 gas in three fresh containers.
Table~\ref{tab:perfrace-cv-distribution}(a) shows that 84 of the 99 wall-clock tasks have a final CV below 6\%.

\textbf{Integrity.}
We validate Integrity across six attack classes: empty or no-op solutions, skipped work, hard-coded outputs, precomputed or cached outputs, modification of protected inputs or tests, and access to verifier-only data.
All 100 accepted tasks completed an explicit red-team pass before Integrity approval.
Table~\ref{tab:perfrace-cv-distribution}(b) shows that 40 tasks were tested against at least two red-team cases, including two tasks with seven or eight cases.
For example, SCC is a source-code statistics tool whose task scans the Linux v6.19 tree and emits JSON statistics; during red-team testing, the agent replaced the scanner with a program that returned the precomputed Linux result instead of scanning the source tree, which passed the workload available to the agent but failed when the verifier asked it to scan a separate pinned checkout of SCC's own source code.
\begin{table}[H]
  \centering
  \small
  \setlength{\tabcolsep}{7pt}
  \renewcommand{\arraystretch}{1.15}
  \begin{minipage}[t]{0.34\textwidth}
    \vspace{0pt}
    \centering
    \textbf{(a) Final measurement CVs}\par\smallskip
    \begin{tabular}{@{}lr@{}}
      \toprule
      CV range & Tasks \\
      \midrule
      \(0\% \leq \mathrm{CV} < 2\%\) & 20 \\
      \(2\% \leq \mathrm{CV} < 6\%\) & 64 \\
      \(6\% \leq \mathrm{CV} < 7\%\) & 8 \\
      \(7\% \leq \mathrm{CV} < 8\%\) & 3 \\
      \(8\% \leq \mathrm{CV} < 9\%\) & 4 \\
      \(9\% \leq \mathrm{CV} < 10\%\) & 0 \\
      \bottomrule
    \end{tabular}
  \end{minipage}
  \hspace{0.035\textwidth}
  \begin{minipage}[t]{0.27\textwidth}
    \vspace{0pt}
    \centering
    \textbf{(b) Red-team tests}\par\smallskip
    \begin{tabular}{@{}rr@{}}
      \toprule
      Tests per task & Tasks \\
      \midrule
      1 & 60 \\
      2 & 20 \\
      3 & 12 \\
      4 & 6 \\
      7 & 1 \\
      8 & 1 \\
      \bottomrule
    \end{tabular}
  \end{minipage}
  \caption{Reliability and integrity evidence: final CVs for the 99 wall-clock tasks and red-team tests for all 100 tasks, with one test assigned conservatively when an exact count cannot be recovered.}
  \label{tab:perfrace-cv-distribution}
\end{table}


\subsection{Evaluation Protocol}
\label{app:perfrace-evaluation}

For each submission, the Judge snapshots the agent's current container and starts a temporary shadow container from that snapshot.
The Judge creates the snapshot with \texttt{docker compose commit}, launches the shadow container under a separate Compose project, and removes its container, local image, and volumes after verification.
The agent container remains paused during verification and resumes from the same workspace state after the shadow container is removed.

Inside the shadow container, the verifier runs through \texttt{tests/test.sh} and emits the correctness verdict and performance metric as JSON.
The verifier code, workload copy, held-out cases, and expected outputs are absent from the agent container and available only inside the shadow container.
Each task declares its warmups, repetition count, aggregation rule, and required correctness pass rate in \texttt{task.toml}, and its verifier returns a performance metric only when that pass rate is met.

The task clock continues while the shadow is created, executed, and removed, so every submission trades part of budget \(B\) for additional correctness and performance feedback.
After verification, the Judge returns the correctness verdict and performance metric to the agent as feedback for subsequent optimization.
For a correct submission, the Judge computes the speedup score by dividing the unmodified baseline metric by the submission metric.
Submissions that fail correctness, exceed the verifier timeout, or do not produce a valid performance metric cannot update the best-so-far trajectory \(S(t)\).

Five minutes before the deadline, the agent receives a final delivery prompt asking it to leave the project in a coherent state before the Judge performs one final shadow verification of the current workspace.
Figure~\ref{fig:perfrace-verification-time} characterizes the time required for one complete Judge evaluation, showing that baseline Judge evaluation remains within the 20-minute limit for all 100 final tasks.
\begin{figure}[H]
  \centering
  \includegraphics[width=0.82\textwidth]{figures/perfrace-verification-time.pdf}
  \caption{Distribution of complete baseline Judge evaluation time across the 100 final task packages, with a mean of 6.1 minutes, a median of 4.9 minutes, and a 20-minute evaluation limit.}
  \label{fig:perfrace-verification-time}
\end{figure}

\phantomsection
\textbf{Longer optimization budgets.}
\label{app:perfrace-longer-budgets}
To measure the benefit of additional optimization time, we run Terminus-KIRA with GPT-5.6-Luna at high reasoning effort three times on each of the 88 evaluated tasks, maintaining one persistent agent session across Judge submissions during every four-hour run.

The aggregate best-so-far correct speedup rises from \(1.314\times\) after one hour to \(1.377\times\), \(1.396\times\), and \(1.435\times\) after two, three, and four hours, respectively.
Across the resulting 264 trajectories, Figure~\ref{fig:perfrace-late-budget-gains} shows that extending
Terminus-KIRA's budget from two to four hours improves the overall score by
4.2\%, and the fourth hour alone adds 2.8\%.

The four-hour score is at least 5\% higher on 24 tasks than at two hours and on
12 tasks than at three hours, showing that agents continue to find meaningful
improvements late in the task.
Among the 228 runs that improve over baseline, the best result is first reached at a median of 88.1\% of the four-hour budget, and 51 of the 264 runs produce a new improvement of at least 5\% during the final hour.

\begin{figure}[H]
  \centering
  \includegraphics[width=\textwidth]{figures/perfrace-long-horizon.pdf}
  \caption{Performance improvements from additional optimization time, based on
  three four-hour runs for each of 88 \perfrace{} tasks. Left: increase in the overall
  score at four hours relative to two and three hours; horizontal lines show
  95\% bootstrap intervals. Right: number of tasks in each improvement range.}
  \label{fig:perfrace-late-budget-gains}
\end{figure}
